\section{总结与延伸}

本讲把语言模型历史、无监督学习、跨模态生成与 code evaluation 串成一条方法论：

\begin{enumerate}
  \item \textbf{生成质量是分层的}：词汇、句法、语义、事实和任务成功不能混成“像不像”；
  \item \textbf{Perplexity 只测分布预测}：它不自动证明真实性、可控性或 downstream utility；
  \item \textbf{Autoregressive objective 提供密集自监督}：analysis by synthesis 可产生可迁移 representation，但也可能学习 shortcut；
  \item \textbf{GPT-1/2/3 改变任务接口}：fine-tuning、zero-shot continuation 和 in-context demonstrations 依次扩大复用；
  \item \textbf{跨模态需要重新定义 token 与 metric}：iGPT、DALL-E 证明 recipe 可迁移，不证明领域结构不重要；
  \item \textbf{Code 需要 functional verifier}：字符串 overlap 不能替代执行语义；
  \item \textbf{Pass@k 测分布中的 rare success}：它与 pass@1 对应不同 candidate budget 和产品体验；
  \item \textbf{Temperature 应按 k 校准}：低温追求单样本质量，高温可在 verifier 支持下增加搜索覆盖；
  \item \textbf{Sampling 只能放大已有概率质量}：若正确程序概率为零，再多 sample 也无用；
  \item \textbf{Production reliability 需要完整流程}：sandbox、tests、static analysis、review、security 与 maintenance 不能由 benchmark 代替。
\end{enumerate}

\begin{importantbox}{Generative-System Evaluation 作业}
选文本、图像或代码任务：定义 model distribution、decoding policy、candidate budget、verifier 与 final selector；同时报告 pass@1/quality、pass@k/diversity、cost、latency 和 failure taxonomy。对代码任务推导 pass@k estimator，并设计三类隐藏 tests；对文本/图像任务说明为什么你的 verifier 不能成为完美 oracle。
\end{importantbox}

\begin{knowledgebox}{最终自检：三个“存在”必须分开}
\textbf{存在于分布}：某个正确样本有非零概率；\textbf{能被找到}：当前 temperature、$k$ 和 search 能采到；\textbf{能被识别}：tests/ranker 能可靠选出。生成模型产品的多数差异，都发生在这三层之间。
\end{knowledgebox}

\subsection{拓展阅读}

{\small
\begin{itemize}[nosep,leftmargin=*]
  \item \href{https://www.cs.utoronto.ca/~ilya/pubs/2011/LANG-RNN.pdf}{\textbf{Generating Text with Recurrent Neural Networks}}：RNN language generation。
  \item \href{https://arxiv.org/abs/1602.02410}{\textbf{Exploring the Limits of Language Modeling}}：large LSTM language modeling。
  \item \href{https://arxiv.org/abs/1704.01444}{\textbf{Learning to Generate Reviews and Discovering Sentiment}}：sentiment neuron 与 unsupervised representation。
  \item \href{https://cdn.openai.com/research-covers/language-unsupervised/language_understanding_paper.pdf}{\textbf{GPT-1}} 与 \href{https://cdn.openai.com/better-language-models/language_models_are_unsupervised_multitask_learners.pdf}{\textbf{GPT-2}}：pretrain/fine-tune 到 zero-shot。
  \item \href{https://arxiv.org/abs/2005.14165}{\textbf{GPT-3}}：few-shot/in-context evaluation 与 human detection study。
  \item \href{https://arxiv.org/abs/2006.10738}{\textbf{iGPT}}：generative pretraining from pixels。
  \item \href{https://arxiv.org/abs/2102.12092}{\textbf{DALL-E}}：zero-shot text-to-image generation。
  \item \href{https://arxiv.org/abs/2107.03374}{\textbf{Codex}}：HumanEval、pass@k 与 code generation evaluation。
\end{itemize}
}

\end{document}
